%% Material moved out of \Cref{sec:setting} to keep the main text short.  Included from the
%% appendix (\Cref{app:shift}); shared by ../paper/main.tex and ../submission/body.tex.

\subsection{Where the contraction condition comes from, and why it fails}
\label{app:contract}

Linearising \eqref{eq:step} at $\ts$ and ignoring $\Pi_t$ gives the error recursion
\begin{equation}
\theta_t-\ts=\big(I-\gamma_t\Hb_{t-1}^{-1}\Hh_t\big)(\theta_{t-1}-\ts)
             -\gamma_t\Hb_{t-1}^{-1}\eps_t+o(\norm{\theta_{t-1}-\ts}).
\label{eq:err}
\end{equation}
Since $\Hb^{-1}\Hh_t$ is a product of two positive semidefinite matrices its eigenvalues are real and
non-negative, so the spectral radius of the linear map in \eqref{eq:err} is below one exactly when the
first inequality of \eqref{eq:contract} holds; the second is sufficient because
$\lambda_{\max}\le\opnorm{\cdot}$. We work with the operator norm because it is what we can compute and
because it is the conservative choice, and we say ``spectral radius'' rather than ``contraction''
deliberately: for a non-symmetric map a spectral radius below one bounds the \emph{asymptotic} growth of
powers, not the norm of a single step. What \eqref{eq:contract} rules out is the regime we actually
observe, in which $\gamma_t\lambda_{\max}\gg2$ and a single step multiplies the error by a large factor.

The per-block maxima of $\opnorm{H^{-1}\Hh_t}$ at $b=20$ are $55$ on the autoregressive design and $89$
on the regime-switching design, so the autoregressive designs are lucky rather than safe. Taking a
full-length step under those numbers makes the first iterates overshoot by orders of magnitude, and
because $\gamma_t=1/t$ undoes an excursion only at rate $1/t$, at realistic $N$ an overshoot in the
first few blocks can survive to the end of the stream. That is the mechanism behind the divergence
reported in \Cref{tab:ablation}(i): without the safeguard, two of six replicates on the
regime-switching design diverge outright, and none with it. This is not a defect of the base method ---
it is created by \emph{blocking} the gradient, which is the one thing we add.

\subsection{What \texorpdfstring{\Cref{prop:shift}}{the safeguard proposition} does not prove}
\label{app:notproved}

The obstruction is real and worth naming: the factor
$\min\{1,c_D(1+\norm{\theta_{t-1}})/\norm{v}\}$ depends on the current block's gradient, so it is
$\mathcal F_t$- rather than $\mathcal F_{t-1}$-measurable and correlates with the noise, which
breaks the martingale structure the proofs use. A complete treatment is available in the
stochastic-approximation literature on expanding truncations
\citep{andrieu2005stability,kushner2003stochastic}, at the price of machinery we do not need here.
\Cref{app:shift} gives the two concrete routes that would close the gap. A unit test asserts the
measured activation counts.

\subsection{The alternative we rejected, and why}
\label{app:rejected}

The contraction condition \eqref{eq:contract} suggests a different
remedy: shift the schedule to $\gamma_t=(t+t_0)^{-1}$ with
$t_0=\tfrac12\max_{t\le\Tw}\opnorm{\Hb_{\Tw}^{-1}\Hh_t}$ measured on the warm-up blocks, which
makes \eqref{eq:contract} hold at $t=1$ by construction and has a cleaner proof (it is a finite
deterministic reindexing; see \Cref{app:shift}). We implemented it and it is worse, for a reason
worth stating: it slows \emph{every} step, whereas the overshoot needs correcting on only a
handful. On the regime-switching design at $N=200{,}000$ the shift leaves the relative variance at
$\ShiftEffMarkovShift$ against $\ShiftEffMarkovCap$ for the safeguard, and on a stream short enough
that $t_0$ is comparable to the number of steps --- which is the case for our real-data segments --- it freezes the estimator
altogether. Both variants are in the code and both are reported
(\Cref{tab:ablation}(i)); the safeguard is the default.



\subsection{Sizing the curvature warm-up}
\label{app:warmup}

It matters: on a $4000$-observation stream with $d=11$, spending $c_{\mathrm w}d=2200$
observations on curvature leaves $\WarmFracStepsUncapped$ blocked steps and a relative variance of
$\WarmFracEffUncapped$, whereas capping at $10\%$ leaves $\WarmFracStepsCapped$ steps and a
relative variance of $\WarmFracEffCapped$ (\Cref{tab:ablation}(j)). Our real-data segments are of exactly that length.

\subsection{Cost accounting for the one-off terms}
\label{app:costdetail}

The warm-up performs $c_{\mathrm w}d$ rank-one updates at $O(d^2)$ each, which is the unavoidable cost
of producing a $d\times d$ curvature estimate at all and is of the same order as a single matrix
factorisation. The step-size shift of \Cref{app:t0} adds a further one-off
$O(c_{\mathrm w}(d^2+db)d/b)$, which is $4$--$5\%$ of the total operation count at the settings of
\Cref{tab:cost} and vanishes relative to the streaming pass as $N$ grows. \Cref{tab:cost} reports the
crossover $N\gtrsim c_{\mathrm w}d^2$ explicitly rather than hiding it in the constant.

\subsection{Dependence in the covariates is not enough}
\label{app:mds}

\begin{remark}[Dependence in the covariates is not enough]
\label{rem:mds}
If the conditional law of $y$ given $x$ is correctly specified and its innovations are
independent over time, then $(g_i)$ is a martingale difference sequence and $S=\Gamma_0$
\emph{exactly}, however persistent $x$ is. A genuine long-run-variance problem needs the
\emph{score} to be serially correlated: serially correlated errors, a latent dynamic
factor, or omitted dynamics. All our designs therefore couple dependent covariates with a
dependent error or latent process, and we report the resulting $\kappa_j$ explicitly.
Covariate dependence still matters for the effective sample size of the curvature
estimate, which is what \Cref{prop:gap} is about.
\end{remark}

\subsection{Misspecification and where dependence has to live}
\label{app:misspec}

In more detail, nothing in the analysis requires the model to be correctly specified: all that is used is that $\ts$ solves
$\nabla F(\ts)=0$ for the \emph{working} loss, so under misspecification $\ts$ is the pseudo-true
parameter and $H^{-1}SH^{-1}$ the corresponding long-run sandwich
\citep{white1982mle,domowitz1982misspecified}. That matters here, because omitted dynamics are one
of the commonest ways for a score to become serially correlated --- and correlated \emph{scores},
not merely correlated covariates, are what the problem requires. If the conditional law of $y$ given
$x$ is correctly specified with innovations independent over time, then $(g_i)$ is a martingale
difference sequence and $S=\Gamma_0$ \emph{exactly}, however persistent $x$ is (\Cref{rem:mds}). All
our designs therefore couple dependent covariates with a dependent error or latent process, and we
report the resulting $\kappa_j$ explicitly.


This matters in practice, because omitted dynamics are one of the commonest ways for a score to become
serially correlated \citep{white1982mle,domowitz1982misspecified}.

\subsection{Why the ridge scale is frozen}
\label{app:freeze}

With $\varsigma_\star$ fixed, positive and finite almost surely under \Cref{as:curv}, \Cref{lem:ridge}
holds unconditionally --- for every $t$, with no assumption on the iterates. The scale exists only to make
$c_\iota$ dimensionless; setting $\varsigma_\star=1$ recovers an unscaled ridge, and
\Cref{tab:ablation}(d) shows the results are the same.
